\documentclass[12pt]{article}

\usepackage{sbc-template}
\usepackage{graphicx,url}
\usepackage[utf8]{inputenc}
\usepackage[brazil]{babel}

\sloppy

\title{Monitoramento e Controle de Tráfego Urbano com Semáforos Inteligentes em um Sistema Distribuído}

\author{
Bruno Braga,
Caio Gomes,
Kaio Henrique Lucio e Santos \\
Lucas Araujo,
Nalberth Vieira,
Pedro Henrique Bellone,
Vitor Alexandre
}

\address{Pontifícia Universidade Católica de Minas Gerais (PUC Minas)\\
Belo Horizonte -- MG -- Brasil
}

\begin{document}

\maketitle

\begin{resumo}
Este trabalho apresenta o desenvolvimento e a evolução de um sistema distribuído para monitoramento e controle de tráfego urbano inteligente e resiliente. A segunda entrega aprofunda a infraestrutura da solução com a migração da camada de transporte para uma Malha HTTP (\textit{HTTP Mesh}) baseada em Express, substituindo o gRPC utilizado na primeira fase. Foram incorporados mecanismos robustos de tratamento de falhas, um algoritmo de escalonamento dinâmico por análise de carga global, transição suave de sinais de trânsito (\textit{preempção segura}) e um sistema de telemetria em tempo real. O painel de controle do Atuador de Borda foi publicado via GitHub Pages, permitindo acesso remoto e validação de consistência entre réplicas. Os resultados obtidos evidenciam a viabilidade prática da abordagem distribuída para aplicações de controle urbano em tempo real.
\end{resumo}

\begin{abstract}
This paper presents the development and evolution of a distributed system for intelligent and resilient urban traffic monitoring and control. The second delivery deepens the solution's infrastructure with a transport layer migration from gRPC to an HTTP Mesh based on Express. Robust fault-handling mechanisms, a dynamic load-based scheduling algorithm, smooth signal preemption, and a real-time telemetry system were incorporated. The Edge Actuator dashboard was deployed via GitHub Pages, enabling remote access and replica consistency validation. Results demonstrate the practical feasibility of the distributed approach for real-time urban control applications.
\end{abstract}

\section{Introdução}

Sistemas de trânsito urbano dependem de decisões rápidas e coordenadas para responder a mudanças no fluxo de veículos. Em arquiteturas centralizadas, a indisponibilidade de um servidor principal pode comprometer toda a operação, gerando atrasos, perda de eventos e baixa resiliência.

Com o objetivo de aplicar conceitos de computação distribuída em um problema prático, este trabalho propõe um sistema de monitoramento e controle de tráfego baseado em múltiplos nós cooperantes. Cada nó é responsável por uma zona da cidade, processa eventos locais de trânsito e compartilha informações com os demais para manter uma visão distribuída do estado do sistema.

Na segunda entrega, a solução foi ampliada com foco em tratamento de falhas, otimizações arquiteturais e validação em cenário real de uso simultâneo. As seções seguintes descrevem em detalhe as mudanças implementadas, os mecanismos de tolerância a falhas e os resultados observados.

\section{Problema Escolhido}

O problema abordado é o monitoramento distribuído de tráfego urbano com atuação sobre semáforos inteligentes. Em cidades com múltiplas regiões interdependentes, o aumento do fluxo em uma via pode impactar rapidamente vias vizinhas, exigindo troca contínua de informações entre diferentes pontos do sistema.

A solução simula esse cenário por meio de zonas monitoradas de forma independente. Cada zona recebe dados de contagem de veículos, identifica situações de tráfego intenso com base em limiares predefinidos e emite alertas para regiões adjacentes, possibilitando reações coordenadas e distribuídas.

Esse problema é adequado ao contexto da disciplina porque pode ser naturalmente decomposto em várias unidades de trabalho executadas em máquinas distintas, exigindo comunicação entre nós, ordenação de eventos e continuidade de operação mesmo diante de falhas.

\section{Arquitetura do Sistema}

A arquitetura foi organizada em três camadas principais: borda, integração e processamento distribuído. Essa separação facilita a coleta de dados, a entrada de eventos no sistema e o processamento descentralizado entre as zonas monitoradas.

\subsection{Camada de Borda}

A camada de borda é composta por uma interface Web desenvolvida com React e Vite, publicada via GitHub Pages. Essa interface simula o painel de um semáforo inteligente e permite o envio de dados em tempo real, como a contagem de veículos observada em determinada região.

Na segunda entrega, o painel foi portado com sucesso para o GitHub Pages, tornando-o acessível de forma remota por qualquer dispositivo. Isso permitiu validar o sistema com controle simultâneo via computador e celular na mesma rede local, comprovando a consistência das réplicas ativas.

\subsection{Camada de Integração}

A camada de integração é uma API REST embutida em cada nó (não há um processo separado). Cada nó expõe o endpoint \texttt{POST /report}, que recebe a contagem de veículos do painel e a propaga ao cluster, e \texttt{GET /intersection-status}, que devolve o estado atual das vias. O CORS é habilitado de forma aberta (\texttt{app.use(cors())}) para que o painel hospedado e dispositivos na rede local acessem a API. Essa camada desacopla o front-end do protocolo interno: a entrada de dados é um POST HTTP simples, enquanto a comunicação entre nós usa endpoints internos próprios.

\subsection{Camada de Processamento Distribuído}

O núcleo da solução é formado por três nós independentes (rodando nas portas lógicas \texttt{50051}, \texttt{50052} e \texttt{50053}), cada um executando em processo isolado. Esses nós se comunicam entre si para compartilhar eventos, detectar condições de congestionamento e coordenar respostas distribuídas.

\subsection{Migração para HTTP Mesh}

Na primeira entrega, a comunicação entre os nós utilizava gRPC com arquivos \texttt{.proto}. Identificou-se, durante a fase de testes, uma limitação severa no subsistema de resolução de nomes do ambiente local (\textit{internal DNS parser}), que impedia o funcionamento correto do gRPC em rede local.

Como solução, a camada de transporte foi substituída por uma \textbf{Malha HTTP (\textit{HTTP Mesh})} baseada em Express. Cada nó expõe endpoints HTTP e troca \textit{payloads} no formato JSON com os demais nós. As propriedades distribuídas foram integralmente preservadas: os processos permanecem isolados, operam de forma desacoplada e continuam se comunicando exclusivamente por passagem de mensagens. O Relógio Lógico de Lamport e o algoritmo Bully foram mantidos e adaptados para esse novo protocolo de transporte.

\subsection{Troca de Mensagens}

O fluxo principal do sistema opera da seguinte forma:
\begin{enumerate}
    \item a interface Web envia a contagem de veículos para um nó via \texttt{POST /report};
    \item o nó atualiza seu estado local e registra o evento com \textit{timestamp} lógico de Lamport;
    \item o nó propaga o valor em \textit{multicast} para os demais nós via \texttt{POST /internal/report-traffic}, que atualizam sua réplica local;
    \item se o limiar de 70 veículos é excedido em alguma via, o coordenador prioriza a via de maior carga e ajusta os semáforos;
    \item cada nó verifica periodicamente o coordenador via heartbeat (\texttt{GET /internal/ping}); se ele cair, os nós iniciam a eleição via Bully.
\end{enumerate}

\section{Requisitos Implementados}

\subsection{Relógios Lógicos de Lamport}

Como os nós operam de forma distribuída, não é seguro depender apenas do horário físico de cada máquina para determinar a ordem dos acontecimentos. Para resolver esse problema, o sistema utiliza relógios lógicos de Lamport \cite{lamport}.

Cada evento relevante incrementa o contador lógico local, e esse valor é anexado às mensagens trocadas entre os nós. Quando uma mensagem é recebida, o nó ajusta seu relógio lógico de acordo com a regra de Lamport, preservando uma ordenação consistente dos eventos observados.

Esse mecanismo mostrou-se especialmente importante durante os testes com controle simultâneo via computador e celular, garantindo o correto ordenamento dos eventos mesmo com mensagens chegando de fontes distintas.

\subsection{Eleição de Líder com Bully Algorithm}

Para evitar ponto único de falha, o sistema possui um mecanismo de eleição de líder baseado no algoritmo Bully. Sempre que o nó coordenador deixa de responder, os demais nós iniciam o processo de eleição.

Nesse processo, nós com identificadores mais altos têm prioridade para assumir a coordenação. Ao final, um novo líder é anunciado ao cluster, permitindo que o sistema continue operando após a falha do coordenador anterior. Este mecanismo foi mantido e adaptado para funcionar sobre o protocolo HTTP Mesh.

\subsection{Controle de Réplicas com Consistência Eventual}

Como requisito adicional da segunda entrega (Item II), cada nó mantém em memória uma réplica completa do estado do cruzamento (as quatro vias, em \texttt{intersectionState}). Quando o painel ou um sensor injeta um dado em um nó, ele atua como coordenador momentâneo da escrita: aplica a regra de negócio e propaga o valor em \textit{multicast} para os demais nós via \texttt{POST /internal/report-traffic}. O relógio de Lamport ordena os eventos e as réplicas convergem para o mesmo estado (consistência eventual).

\section{Tratamento de Falhas}

O tratamento de falhas representa o requisito central da segunda entrega e foi implementado em múltiplos níveis do sistema.

\subsection{Detecção de Falha por Heartbeat}

A cada 2 segundos (\texttt{HEARTBEAT\_INTERVAL}), cada nó envia um \texttt{GET /internal/ping} ao coordenador atual. Se o coordenador não responde, o nó o considera indisponível e dispara a reeleição. Falhas de nós comuns são absorvidas pela degradação graciosa: todas as chamadas entre nós ficam dentro de blocos \texttt{try/catch}, então um nó que cai simplesmente deixa de receber as propagações, sem interromper os demais.

\subsection{Reeleição Automática}

Quando o heartbeat detecta a queda do coordenador, o nó chama \texttt{startElection()} automaticamente, sem intervenção manual. O algoritmo Bully elege o nó de maior identificador entre os ativos, que assume a coordenação e avisa os demais via \texttt{POST /internal/set-coordinator}. Nos testes com três nós, a queda do coordenador (Nó 3) foi detectada pelo Nó 1 e o Nó 2 assumiu como novo coordenador automaticamente.

\subsection{Recuperação de Estado na Reconexão}

Para que um nó que reinicia não volte com o estado zerado, ao subir ele chama \texttt{syncStateFromPeers()}: faz um \texttt{GET /internal/state} no primeiro peer ativo que encontrar e adota o estado das quatro vias e o relógio de Lamport recebidos. Durante a operação normal, o estado de tráfego é mantido replicado pela propagação em \textit{multicast} descrita anteriormente.

\subsection{Testes de Falha Realizados}

Foram realizados os seguintes testes para validar o tratamento de falhas:
\begin{itemize}
    \item \textbf{Queda do coordenador}: com os três nós no ar e o Nó 3 eleito coordenador, o processo do Nó 3 foi encerrado. O Nó 1 detectou a ausência de resposta pelo heartbeat e disparou a reeleição; o Nó 2 assumiu como novo coordenador. O failover completo (detecção mais eleição) ficou abaixo de 5 segundos.
    \item \textbf{Falha de nó comum}: um nó não-coordenador foi encerrado. As propagações para ele passaram a falhar silenciosamente (\texttt{try/catch}) e os nós restantes continuaram operando sem interrupção.
    \item \textbf{Reconexão de nó}: um nó reiniciado recuperou o estado atual via \texttt{GET /internal/state} e voltou a participar do cluster com o mesmo estado de vias dos demais.
\end{itemize}

\section{Otimizações Implementadas na Segunda Entrega}

\subsection{Escalonamento por Análise de Carga Global}

Na primeira versão, o sistema forçava a abertura de uma via de forma abrupta e isolada caso ela ultrapassasse 70 veículos, quebrando o ciclo ordenado de alternância. Na segunda entrega, foi implementado um algoritmo de \textbf{Alocação Dinâmica baseada em Carga Máxima}.

O Coordenador avalia continuamente o vetor de tráfego de todas as quatro vias. Quando um alerta crítico é disparado, o sistema compara os pesos de todas elas e prioriza estritamente aquela com o maior congestionamento real, em vez de atender simplesmente o primeiro alerta recebido. Isso evita que vias com carga intermediária monopolizem o acesso verde em detrimento de vias em situação mais crítica.

\subsection{Preempção Segura de Sinais}

Na versão anterior, o semáforo chaveava instantaneamente de Verde para Vermelho em caso de prioridade, comportamento considerado perigoso em um cenário real. Na segunda entrega, foi desenvolvido um estado de \textbf{Preempção Suave} (\texttt{isTransitioning}).

Quando uma via crítica solicita passagem, a via atualmente aberta entra obrigatoriamente em estágio Amarelo por 3 segundos, permitindo o escoamento seguro dos veículos antes de fechar e ceder o Verde à via prioritária. Esse mecanismo torna a simulação mais realista e segura.

\subsection{Simulação de Vazão Contínua}

A função de geração aleatória de tráfego a cada 12 segundos foi removida, pois prejudicava os testes manuais e mascarava o comportamento real do sistema. O sistema passou a reagir exclusivamente aos dados inseridos pelos sensores ou pelo painel Web.

Além disso, foi implementado um mecanismo de \textbf{Vazão Contínua}: enquanto uma via permanece em Verde, o servidor consome gradativamente a carga de veículos daquela via (reduzindo 2 carros por segundo), simulando o escoamento dos veículos ao longo do tempo.

\subsection{Telemetria em Tempo Real}

Foi implementada uma função centralizada de telemetria no console (\texttt{console.table}). A cada mudança de estado, seja pelo ciclo Round-Robin, pela chegada de dados de sensor ou pela ativação de via crítica, todos os nós imprimem instantaneamente uma tabela formatada no terminal, exibindo o status atual, o tempo restante em segundos e a carga de veículos de cada uma das quatro vias.

Essa funcionalidade facilitou significativamente a identificação de comportamentos inesperados durante os testes e a depuração de problemas de sincronização.

\section{Análise de Desempenho}

Para avaliar o impacto da distribuição, medimos o tempo de uma escrita (\texttt{POST /report}) no modo distribuído (o Nó 1 processa e propaga via HTTP Mesh para os outros dois nós) e no modo centralizado de referência (a mesma operação feita só em memória, sem rede). A medição foi feita com o script \texttt{benchmark.js}, com 50 escritas por modo, na mesma máquina. A Tabela~\ref{tab:latencia} resume os tempos obtidos.

\begin{table}[ht]
\centering
\caption{Tempo de uma escrita (\texttt{POST /report}) -- 50 medições por modo.}
\label{tab:latencia}
\begin{tabular}{lccc}
\hline
\textbf{Modo} & \textbf{Média} & \textbf{Mínimo} & \textbf{Máximo} \\
\hline
Distribuído (HTTP Mesh + propagação) & 6,4 ms & 4,7 ms & 40 ms \\
Centralizado (memória, sem rede) & $<$ 0,1 ms & $\sim$0 ms & 0,02 ms \\
\hline
\end{tabular}
\end{table}

Como esperado, o modo centralizado é praticamente instantâneo por não envolver rede, enquanto o distribuído adiciona um \textit{overhead} de poucos milissegundos por escrita, devido à malha HTTP e à propagação para as réplicas.

O diferencial do modo distribuído aparece na falha, resumido na Tabela~\ref{tab:comparativo}: no modo centralizado a queda do processo único interrompe todo o sistema, enquanto no distribuído a operação é retomada automaticamente em menos de 5 segundos após a queda do coordenador, sem perda do estado de tráfego.

\begin{table}[ht]
\centering
\caption{Comparativo entre os modos centralizado e distribuído.}
\label{tab:comparativo}
\begin{tabular}{lcc}
\hline
\textbf{Critério} & \textbf{Centralizado} & \textbf{Distribuído} \\
\hline
Latência por escrita & $<$ 0,1 ms & $\sim$6,4 ms \\
Ponto único de falha & Sim & Não \\
Comportamento na queda & Sistema para & Reeleição automática \\
Tempo de recuperação & Manual & $<$ 5 s \\
Perda de estado na falha & Total & Nenhuma \\
\hline
\end{tabular}
\end{table}

Para o problema de tráfego urbano, em que as decisões ocorrem em escala de segundos, um custo de poucos milissegundos por escrita é irrelevante diante do ganho de disponibilidade.

\section{Desafios e Problemas Encontrados}

O principal desafio desta segunda entrega foi a migração da camada de transporte de gRPC para HTTP Mesh. A limitação do subsistema de DNS interno do ambiente de desenvolvimento local impedia a resolução correta dos endereços necessários para o funcionamento do gRPC, exigindo uma revisão completa da camada de comunicação. A migração demandou a adaptação de todos os pontos de comunicação entre nós e a revalidação dos mecanismos de Lamport e Bully sobre o novo protocolo.

Outro desafio foi manter o sistema operando durante a reeleição do líder. Como cada nó mantém sua própria réplica do estado e roda o laço dos semáforos localmente, a troca de coordenador não interrompe o controle do cruzamento; um nó que reinicia recupera o estado via \texttt{GET /internal/state}. Com isso, evitamos uma janela em que o cruzamento ficaria sem controle durante a eleição.

A publicação do painel no GitHub Pages também exigiu habilitar o CORS nos nós (\texttt{app.use(cors())}), permitindo que o painel hospedado e dispositivos na rede local acessem a API HTTP dos nós.

\section{Papel dos Integrantes}

A seguir, apresenta-se um resumo atualizado do papel de cada integrante no desenvolvimento do trabalho, incluindo as contribuições específicas da segunda entrega.

\begin{itemize}
    \item \textbf{Kaio Henrique Lucio e Santos}: coordenação geral do projeto, definição e revisão da arquitetura distribuída, migração da camada de transporte para HTTP Mesh e integração entre os componentes.
    \item \textbf{Bruno Braga}: desenvolvimento e manutenção da interface Web (React + Vite) e publicação do painel no GitHub Pages.
    \item \textbf{Caio Gomes}: implementação da API REST de integração embutida nos nós (\texttt{/report} e \texttt{/intersection-status}), configuração do CORS e integração com o núcleo distribuído.
    \item \textbf{Lucas Araujo}: implementação do algoritmo de Alocação Dinâmica baseada em Carga Máxima e do mecanismo de Preempção Suave de sinais.
    \item \textbf{Nalberth Vieira}: adaptação dos relógios lógicos de Lamport para o protocolo HTTP Mesh e testes de ordenação de eventos.
    \item \textbf{Pedro Henrique Bellone}: desenvolvimento do tratamento de falhas, incluindo detecção por \textit{heartbeat}, reeleição automática via Bully e sincronização de estado pós-falha.
    \item \textbf{Vitor Alexandre}: testes de falha e reconexão, implementação do sistema de telemetria em tempo real e documentação técnica do trabalho.
\end{itemize}

\section{Conclusão}

O trabalho desenvolvido aplica e aprofunda conceitos fundamentais de sistemas distribuídos em um cenário realista de monitoramento e controle de tráfego urbano. A segunda entrega avançou significativamente sobre a base construída anteriormente, incorporando tratamento robusto de falhas, otimizações no algoritmo de escalonamento, transição segura de sinais e uma infraestrutura de borda acessível remotamente.

A migração do gRPC para o HTTP Mesh, embora motivada por uma limitação de ambiente, demonstrou que as propriedades distribuídas fundamentais, isolamento de processos, passagem de mensagens, ordenação lógica de eventos e eleição de líder, são independentes do protocolo de transporte subjacente, desde que a comunicação permaneça desacoplada e assíncrona.

Como possíveis melhorias futuras, destacam-se: a implementação do protocolo Two-Phase Commit (2PC) para garantia de atomicidade em transações que afetam múltiplas vias simultaneamente; a expansão do sistema para um número maior de zonas e cruzamentos interligados; e a integração com dados reais de sensores de tráfego via APIs públicas de mobilidade urbana.

\bibliographystyle{sbc}
\begin{thebibliography}{99}

\bibitem{lamport}
LAMPORT, Leslie. ``Time, Clocks, and the Ordering of Events in a Distributed System.'' \textit{Communications of the ACM}, v. 21, n. 7, p. 558--565, 1978.

\bibitem{tanenbaum}
TANENBAUM, Andrew S.; VAN STEEN, Maarten. \textit{Distributed Systems: Principles and Paradigms}. 2. ed. Upper Saddle River: Prentice Hall, 2007.

\bibitem{garcia-molina}
GARCIA-MOLINA, Hector. ``Elections in a Distributed Computing System.'' \textit{IEEE Transactions on Computers}, v. C-31, n. 1, p. 48--59, 1982.

\bibitem{express}
OpenJS Foundation. ``Express: Fast, unopinionated, minimalist web framework for Node.js.'' Disponível em: \url{https://expressjs.com/}. Acesso em: jun. 2026.

\bibitem{react}
Meta Open Source. ``React: A JavaScript library for building user interfaces.'' Disponível em: \url{https://react.dev/}. Acesso em: jun. 2026.

\end{thebibliography}

\end{document}